\section{Auxiliary structural reductions}\label{app:structural-auxiliary}
The following arguments are useful independently of the width-two proof.
They concern the stated local laws or reductions and make no claim about
arbitrary larger-width extensions.

\begin{lemma}[Canonical pair laws on a forest]\label{lem:canonical-pairs}
On a forest of binary variables of prescribed means $p_i$, let $Q_{ij}$
be the midpoint of the two Fr\'echet endpoint pair laws on each edge:
\[
 Q_{ij}(1,1)=\frac{\min(p_i,p_j)+\max(0,p_i+p_j-1)}2.
\]
For any target set $S$ there is a global law with these pair marginals
and target union probability at least
$[\min(1,\sum_{i\in S}p_i)+\max_{i\in S}p_i]/2$, with both expressions
zero for $S=\varnothing$. For any nonnegative target payoff, such a law
also attains at least half its unconstrained singleton-feasible maximum.
The resulting law may depend on the target.
\end{lemma}
\begin{proof}
Every pair cell is affine in its $(1,1)$ probability and nonnegative
on the Fr\'echet interval. Its midpoint value is therefore at least
half every feasible cell value. Choose a singleton-feasible law $P$
maximizing the target union, using consecutive circular event intervals
to attain $c=\min(1,\sum_Sp_i)$. The edge measures
$R_{ij}=2Q_{ij}-P_{ij}$ are nonnegative probability measures with the
prescribed singleton means. They glue on the forest by rooting each
component and sampling each child conditionally on its parent. Null
conditioning states are irrelevant. Thus a law $R$ with these edge
marginals exists. The mixture $(P+R)/2$ has exactly the desired $Q$
marginals, and its target union is at least $(c+b)/2$ since every law
has union at least $b=\max_Sp_i$. For a nonnegative payoff use its
maximizing law $P$ and $\E_Rg\ge0$ in the same construction.
\end{proof}
This residual construction cannot simply be applied on a triangle:
for three fair bits, take $P$ equally on the two constant states.
Each canonical pair is uniform, so $2Q_{ij}-P_{ij}$ requires perfect
anticorrelation. Three bits cannot be pairwise anticorrelated. This
is a failure of that residual-gluing step, not a counterexample to a
positive-monomial width-two bound.

\begin{lemma}[Identical-neighborhood compression]\label{lem:twin-compression}
In a positive unit-cube monomial factorization, suppose a nonempty
variable block $B$ occurs in every incident factor in its entirety.
Replace it by one bit $W$ of mean
$w=\max(0,\sum_{i\in B}x_i-|B|+1)$ in those factors.
The full lower envelope and sum of local lower envelopes equal those
of the reduced instance. For some $\Delta\ge0$,
\[
 \tgap_{\rm original}=\tgap_{\rm reduced}+\Delta,\qquad
 \hgap_{\rm original}=\hgap_{\rm reduced}+\Delta.
\]
Thus every reduced gap inequality with constant at least one transfers.
\end{lemma}
\begin{proof}
For any original binary law, the product $V=\prod_{i\in B}X_i$ has
mean $q\ge w$. If $q>0$, independently retain its ones with probability
$w/q$, producing $W\le V$ of mean $w$; if $q=0$ use $W=0$.
The reduced objective is nondecreasing in $W$, so this operation cannot
increase it. Conversely choose a local law on $B$ attaining product
mean $w$ with all the original means. Attach it to a reduced law by
conditioning on its product value $W$. This lifts the law with the
same objective and original means. These two operations prove full
lower equality, including the endpoint cases.
For a factor with scope $B\cup C$, its original local lower value is
\[
 \max(0,\sum_Bx_i+\sum_Cx_i-|B|-|C|+1)
 =\max(0,w+\sum_Cx_i-|C|),
\]
where if $\sum_Bx_i-|B|+1<0$ both sides vanish. The upper value can
only decrease under compression because $w\le\min_Bx_i$; this also
holds when $C$ is empty and the reduced factor is affine.
Common upper attainment permits summing these decreases into one
$\Delta\ge0$. The two gap identities follow, and
$\tgap_{\rm reduced}+\Delta\le c(\hgap_{\rm reduced}+\Delta)$ for
$c\ge1$ proves the transfer without dividing by a gap.
\end{proof}

\begin{proposition}[Signed forest gluing on arbitrary boxes]
\label{prop:signed-forest}
For any multiaffine local factors on finite boxes whose incidence graph
is a forest, both upper and lower envelopes add. In particular this
holds for frequency-two factors whose dual multigraph is a forest.
\end{proposition}
\begin{proof}
Choose an optimizing vertex law for each factor and normalize its
coordinate endpoints. Root each incidence-tree component at a factor
and traverse it, sampling each new factor conditionally on its single
already sampled variable. Consistent singleton means ensure consistency
at that separator; every other variable is new, or there would be an
incidence cycle. Null separator states have zero probability. This
simultaneously preserves all chosen optimizing laws. Use minimizing
laws for the lower identity and maximizing laws for the upper identity.
Different components and unused variables can be sampled independently.
Fixed coordinates are substituted first. Signs do not enter the proof.
\end{proof}

A factor partition into $q$ incidence forests would give a $q$ bound
by mixing their exact laws for positive monomials, but treewidth does
not bound such a partition number. In $K_{2,m}$, $m\ge2$, with the two-vertex
side interpreted as variables, any two factors create a four-cycle;
all $m$ factors need distinct forest colors, although treewidth is two.
Adding one private variable to each factor makes their monomial scopes
distinct and preserves the obstruction. This explains why the stronger
matrix argument is useful.

For the radix construction, the general cutoff following
Theorem~\ref{thm:variable-radix} also records a fixed-radix limitation.
For integers $b\ge2,L\ge2$ its profile capacities satisfy
\[
 M_s=\frac{(L-s)(b-1)+b}{b^s},\qquad
 \sum_{q=s+1}^L M_q=\frac{L-s}{b^s}.
\]
The sum identity follows by telescoping
$M_q=(L-q+1)b^{1-q}-(L-q)b^{-q}$.
If $s$ brackets one between $M_s$ and $M_{s+1}$, the exact hull is
$s+(L-s)b^{-s}$, as proved by its two attaining profiles.
For fixed $b$, $s=\log_b L+O_b(1)$: the inequality $M_s\ge1$ gives
$s\le\log_b((b-1)L+b)$, hence $L-s\sim L$; the reverse bracket then
gives $b^{s+1}\ge(b-1)(L-s-1)+b$. The remainder term is $O_b(1)$,
so the hull is $\log_b L+O_b(1)$. This is a limitation of this
fixed-radix construction, not a sharp second-order lower theorem for
all positive polynomials.
